\clearpage
\setcounter{page}{1}
\setcounter{section}{0}
\renewcommand{\thesection}{\Alph{section}} 
\maketitlesupplementary


\section*{Outline}
\label{sec:outline}
This supplementary material provides two sets of additional information on AnyTalker.
\subsection*{A. Experimental Details}
\begin{itemize}
\item Data-processing pipeline used during training
\item Construction scheme for inference data
\item Training hyper-parameters and other technical details
\item Inference settings
\end{itemize}
\subsection*{B. Extended Experiments}
\begin{itemize}
\item More analysis of the Interactivity metric
\item Additional experimental results
\item Effectiveness of interactivity refinement
\end{itemize}

\section*{A. Experimental Details}
\label{sec:data-processing}
\subsection*{A.1. Training Data Processing}
\label{subsec:training-data-processing}
\begin{figure}[t]
    \centering
\includegraphics[width=\linewidth]{fig/sync.pdf}
    \caption{SyncNet Score Matrix in two-person data.}
    \label{fig:sync}
\end{figure}

\noindent \textbf{Auxiliary Models.}
We first introduce the auxiliary models used during the data processing stage.
To segment sentences and chunk the audio track extracted from each video, we employ the pre-trained speaker-diarization-3.1 model~\cite{Plaquet23}\footnote{\url{https://huggingface.co/pyannote/speaker-diarization-3.1}} in conjunction with OpenAI’s Whisper~\cite{radford2023robust}\footnote{\url{https://github.com/openai/whisper}}.
For vocal separation and vocal feature extraction, we use pre-trained models
from Kim\_Vocal\_2~\cite{takahashi2017multi}\footnote{\url{https://github.com/Anjok07/ultimatevocalremovergui}}
and Wav2Vec2~\cite{schneider2019wav2vec,baevski2020wav2vec}\footnote{\url{https://huggingface.co/facebook/wav2vec2-base-960h}}.
For face detection and facial feature processing, we adopt RetinaFace~\cite{deng2020retinaface}
and ArcFace~\cite{deng2019arcface}, both accessible via InsightFace\footnote{\url{https://github.com/deepinsight/insightface}}.
To measure audio-visual synchronization, we employ the pre-trained SyncNet model from LatentSync~\cite{li2024latentsync}\footnote{\url{https://huggingface.co/ByteDance/LatentSync/tree/main}}.
Text prompts are obtained with Gemini~2.5~Pro\footnote{\url{https://deepmind.google/models/gemini/pro/}},
and their features are extracted using the T5-encoder~\cite{raffel2020exploring}\footnote{\url{https://huggingface.co/Wan-AI/Wan2.1-I2V-14B-720P/tree/main/google/umt5-xxl}}.
Finally, we filter videos with excessive camera motion by means of the CoTracker model~\cite{karaev2024cotracker}\footnote{\url{https://huggingface.co/facebook/cotracker}}.

\begin{figure}[t]
    \centering
\includegraphics[width=0.7\linewidth]{fig/data.pdf}
    \caption{Properties of training data.}
    \label{fig:data}
\end{figure}

\begin{figure*}[t]
    \centering
\includegraphics[width=\linewidth]{fig/multi-data.pdf}
    \caption{
    Two cases from InteractiveEyes. 
    Both of the speakers have speaking periods.
    }
    \label{fig:multi-data}
\end{figure*}

\begin{figure}[t]
    \centering
\includegraphics[width=\linewidth]{fig/single-data.pdf}
    \caption{Input cases from the benchmark dataset, from left to right: HDTF~\cite{zhang2021flow}, VFHQ~\cite{xie2022vfhq}, and EMTD~\cite{meng2025echomimicv2}.
    }
    \label{fig:single-data}
\end{figure}

\noindent \textbf{Single-Person Data.}
A valid single-person clip is a 24-fps video that continuously exhibits one identifiable face and is accompanied by speech audio whose lip movements are perfectly synchronised with the soundtrack.
We first eliminate hand-held recordings exhibiting pronounced camera shake or abrupt scene changes using CoTracker~\cite{karaev2024cotracker}.
RetinaFace~\cite{deng2020retinaface} is then applied to guarantee that most frame contains exactly one face.
Videos that satisfy the above criteria are processed by speaker-diarization-3.1~\cite{Plaquet23} to obtain coarse sentence-level segments.
Utterances longer than 5\,s are further subdivided into semantically coherent units with Whisper~\cite{radford2023robust}, after which any fragment containing overlapping speakers is discarded.
The remaining clean single-speaker segments average $\approx$\,2\,s in duration.
Because 2\,s clips cannot saturate GPU memory, we randomly concatenate consecutive segments to extend each sequence; the resulting clip lengths follow a Gaussian distribution with mean 4\,s and standard deviation 0.5\,s.
Then, their vocal tracks are extracted with Kim\_Vocal\_2~\cite{takahashi2017multi} and encoded with Wav2Vec-2~\cite{schneider2019wav2vec}, while audio-visual synchrony is scored by SyncNet~\cite{li2024latentsync}.
Clips whose synchrony score lies below a pre-defined threshold are removed from the training set.
Finally, every retained clip is transcribed and annotated with frame-level face bounding boxes.
We then use Gemini 2.5 Pro to generate textual descriptions for these clips and employ a T5 encoder to extract text features.
All input text is uniformly truncated or padded to exactly 512 tokens.
After being preceded by the above pipeline, it yields approximately $1,000$ h of high-quality 480P single-person data.


\noindent \textbf{Two-Person Data.}
Processing Two-person clips follows the single-person pipeline with three key modifications.
\begin{enumerate}[leftmargin=*,nosep]
  \item \textbf{Face count:} the video must contain \emph{exactly two} faces in most frame.
  \item \textbf{Speaker activity:} speaker-diarization-3.1~\cite{Plaquet23} is constrained to output only two valid states: (i) both speakers active or (ii) a single speaker active. 
  \item \textbf{Spatial consistency:} the left--right spatial ordering of the two faces must remain unchanged throughout the entire clip; identity swapping is detected and rejected via InsightFace~\cite{deng2020retinaface}.
\end{enumerate}
To establish the correct voice--face correspondence, we compute a $2\!\times\!2$ SyncNet confidence matrix (\cref{fig:sync}) and requires the two largest scores to lie on the diagonal; otherwise, the clip is discarded.
A minimum synchrony threshold identical to the single-person setting is further applied.
After filtering, approximately $12$ h of clean two-person data are retained.

\noindent \textbf{Data Concatenation.}
During the first-stage training, we horizontally concatenate randomly selected single-person clips.  
Because the original videos are extremely high-resolution (many 2K/4K), naively resizing them along the height axis and stacking would yield faces that occupy only a tiny fraction of the frame.  
To avoid this, we adopt a special cropping strategy.  
First, we locate the face centre in each of the two candidate clips.  
For 480P footage, the frame size is (H, W)=(480,832), so we expand a minimal crop of (480,416) around each face centre.  
We then apply further augmentation: while keeping the crop inside the original image, we randomly enlarge the window while preserving the 480/416 aspect ratio.  
The resulting crops contain a much larger facial region, enabling the model to learn a more accurate lip-to-audio mapping.

\noindent \textbf{Data Properties.}
Each processed sample is summarised by a dictionary-style item that stores the absolute paths to the decoded video, cleaned audio, and pre-extracted features.
During training, the data loader indexes all information through this item.
Single-person and two-person clips share an identical key structure; the number of face--speech entries in the list unambiguously indicates whether the sample contains one or two speakers.
Further details are illustrated in~\cref{fig:data}.

\begin{figure*}[t]
    \centering
\includegraphics[width=0.9\linewidth]{fig/case.pdf}
    \caption{
    Four typical cases that will get a high Interactivity score.
    }
    \label{fig:case}
\end{figure*}

\subsection*{A.2. Benchmark Data Processing}
\label{subsec:benchmark-data-processing}

\noindent \textbf{Single-Person Data.}
We adopt HDTF~\cite{zhang2021flow} and VFHQ~\cite{xie2022vfhq} as the single-speaker evaluation benchmarks.
Both datasets are de facto standards for talking-head generation; their images are tightly cropped around the face region.
Additional experiments on half-body sequences that include hands are reported in Sec.~B.2.

For each test video, we supply (i) a single reference image of the subject and (ii) the corresponding speech segment.
To respect GPU-memory constraints, we fix the audio length to 6\,s---a duration that all competing methods process without overflow.
Twenty clips are randomly selected from each dataset, ensuring that none of the identities appear in the AnyTalker training set.
Reference images from the three evaluation splits are visualised in~\cref{fig:single-data}.


\noindent \textbf{Two-Person Data.}
We provide additional details for \emph{InteractiveEyes}, the test set introduced in Section 4 in the main text for two-speaker scenarios.
Each video is shot with a stable camera and contains \emph{exactly two faces in every frame}; the duration is approximately 10\,s.
In 80\,\% of the cases, both speakers produce speech, while in the remaining 20\,\% only one speaker talks and the other maintains a listening pose.
To avoid segmentation errors, we do \emph{not} apply speaker-diarization-3.1~\cite{Plaquet23}; instead, speaking intervals are manually annotated so that the Interactivity metric proposed in the main paper can be computed accurately.
Two representative cases are shown in~\cref{fig:multi-data}.

\begin{table}[t]
    \centering
    \caption{
       Quantitative result on EMTD~\cite{meng2025echomimicv2} benchmark.
    }
    \small
    \begin{tabular}{cccccc}
        \toprule
        \multirow{2}{*}{\textbf{Method}} & \multicolumn{4}{c}{\textbf{Metrics}} \\
        \cmidrule(lr){2-5}
        & \textbf{Sync-C$\uparrow$} & \textbf{FID} & \textbf{FVD} & \textbf{ID$\uparrow$} \\
        \midrule
        FantasyTalking~\cite{wang2025fantasytalking} & 3.76 & 67.66 & 818.71 & 0.76 \\
        EchoMimic v2~\cite{meng2025echomimicv2} & 6.27 & 63.43 & \underline{671.18} & 0.76 \\
        MultiTalk~\cite{kong2025let} & \underline{8.38} & 64.71 & 787.99 & \textbf{0.79} \\
        \midrule
        AnyTalker-1.3B & 5.83 & \underline{56.01} & 789.59 & 0.74 \\
        AnyTalker-14B & \textbf{8.45} & \textbf{50.61} & \textbf{664.58} & \underline{0.77} \\
        \bottomrule
    \end{tabular}
    \label{tab:EMTD}
\end{table}

\begin{figure*}[t]
    \centering
\includegraphics[width=\linewidth]{fig/bind.pdf}
    \caption{
   A bad case generated by Bind-Your-Avatar~\cite{huang2025bind}.
    }
    \label{fig:bad}
\end{figure*}


\begin{figure*}[t]
    \centering
\includegraphics[width=0.9\linewidth]{fig/qualitative_single.pdf}
    \caption{
   Qualitative results on HDTF~\cite{zhang2021flow} (left) and VFHQ~\cite{xie2022vfhq} (right) benchmark.
   The positions of pronunciation have been highlighted in {\color{red}red} and \underline{underlined}.
    }
    \label{fig:qualitative-single}
\end{figure*}

\subsection*{A.3. Implement Details}
\label{sec:exp-detail}

\noindent \textbf{Training Details.}
To comprehensively evaluate our method, we train two models of different sizes: Wan2.1-1.3B-Inp\footnote{\url{https://huggingface.co/alibaba-pai/Wan2.1-Fun-1.3B-InP}} and Wan2.1-I2V-14B~\cite{wan2025wan}, which served as the foundational video diffusion models for our experiments. 
In all stages, the text~\cite{raffel2020exploring}, audio~\cite{baevski2020wav2vec}, and image~\cite{radford2021learning} encoders, as well as the 3D VAE, remained frozen with their parameters unchanged. 
The DiT main network, including the newly added AFCA Layers, had all its parameters open for training. 
The first stage employs a higher learning rate of \(2 \times 10^{-5}\) for pretraining, while the second stage uses a lower learning rate of \(5 \times 10^{-6}\) for fine-tuning, incorporating a warm-up strategy and optimized using the AdamW optimizer~\cite{loshchilov2017decoupled}. 
The 14B model is trained using 32 NVIDIA H200 GPUs. 
In the first stage, the global batch size is set to 32 and the model is trained for 2.4M steps. 
In the second stage, the batch size is adjusted to 16, and the model is trained for an additional 50K steps.
The 1.3B model is trained using 8 NVIDIA H200 GPUs. 
The global batch size is maintained at 48 throughout the training process, with the total number of training steps being consistent with that of the 14B model.


\noindent \textbf{Inference Details.}
For all competing methods, we strictly follow the publicly released implementations and adopt their default recommended inference hyperparameters.
Approaches that require textual input receive a fixed prompt for the single-person benchmark: ``\texttt{this person is talking}''.
For the multi-person benchmark, the prompts are automatically generated by Gemini~2.5~Pro as described in~Sec.~A.1.
AnyTalker performs inference with classifier-free guidance (CFG)~\citep{ho2021classifier} using a guidance scale of $4.0$; in the unconditional branch, both textual and audio features are set to zero.
Face masks are extracted with InsightFace and uniformly dilated to provide a slightly enlarged region for generation.


\section*{B. Extended Experiments}
\label{sec:exp-res}

\subsection*{B.1. More analysis about Interactivity Metric}
\label{subsec:S-metric}

\begin{figure*}[t]
    \centering
\includegraphics[width=\linewidth]{fig/more.pdf}
    \caption{
   More Results generated by AnyTalker.
    }
    \label{fig:anytalker-result}
\end{figure*}

\begin{figure*}[t]
    \centering
\includegraphics[width=0.6\linewidth]{fig/ab-inter.pdf}
    \caption{
   Improvement in interaction among each identity after fine-tuning on authentic multi-person data (right). 
    }
    \label{fig:ab-inter}
\end{figure*}




\noindent \textbf{Good Case.}
\cref{fig:case} visualises four representative listener behaviours that strongly signal conversational engagement: eyebrow raise, head nod, head turn, and gaze shift.
The presence of such cues contributes substantially to the final Interactivity score.
Because Interactivity is computed over the entire video, we do not report frame-level values in~\cref{fig:case}.
As illustrated in~\cref{fig:anytalker-result}, AnyTalker produces sequences rich in these interactive actions, and thus achieves a comparatively high Interactivity rating.



\noindent \textbf{The Rubustness of Interactive Metric.}
Some generative baselines occasionally produce highly implausible motions.
For example, Bind-Your-Avatar~\citep{huang2025bind} generates an exaggerated ``lying-down'' action shown in~\cref{fig:bad}.
Without counter-measures, such artifacts would dramatically inflate the ${Motion}$ score even though they are unrelated to interactivity.
We therefore introduce a lightweight anomaly-suppression rule: if the mean facial-landmark displacement between two consecutive frames exceeds 10 pixels (all faces are pre-aligned to a $256\times 256$ canvas), the landmark positions are frozen until a subsequent displacement below 10 pixels is observed.
As demonstrated in the right two frames of~\cref{fig:bad}, this simple clamping prevents the vast majority of abnormal movements from entering the Interactivity computation.


\subsection*{B.2. Additional Experimental Results}
\label{subsec:more-quali}

\noindent \textbf{Quantitative Results on the EMTD Benchmark.}
EMTD~\citep{meng2025echomimicv2} is a half-body dataset that includes hands, as illustrated in~\cref{fig:data}.
We compare AnyTalker against three methods capable of generating half-body sequences: EchoMimic\,v2~\citep{meng2025echomimicv2}, FantasyTalking~\citep{wang2025fantasytalking}, and MultiTalk~\citep{kong2025let}.
The evaluation protocol strictly follows the metrics described in Sec.\,5.1 for the single-person benchmark.
AnyTalker-14B achieves the best scores on three metrics and is only marginally behind MultiTalk on identity preservation (ID).
These results demonstrate that AnyTalker is a comprehensive framework that handles not only tight-face inputs but also half-body scenarios.

\noindent \textbf{Qualitative Result on Single-Person Benchmarks.}
As illustrated in~\cref{fig:qualitative-single}, we present qualitative comparisons of EchoMimic~\cite{chen2025echomimic}, StableAvatar~\cite{tu2025stableavatar}, Sonic~\cite{ji2025sonic}, MultiTalk~\cite{kong2025let}, OmniHuman-1.5~\cite{jiang2025omnihuman}, and AnyTalker-14B on the HDTF~\cite{zhang2021flow} and VFHQ~\cite{xie2022vfhq} benchmarks.
AnyTalker consistently produces clear dentition and accurate lip movement.

\noindent \textbf{More Multi-Person Results.}
\cref{fig:anytalker-result} showcases additional multi-person animations generated by AnyTalker. The model gracefully handles a broad spectrum of inputs: real photographs, AIGC images, and cartoons. 
It produces natural, context-appropriate interactions regardless of the number of identities involved. Further compelling examples are available in the videos on the project homepage.

\subsection*{B.3. Effectiveness of Interactivity Refinement}
\label{sec:ab}
As shown in~\cref{fig:ab-inter}, refining interactivity with authentic multi-person data lets identities exchange significantly more natural eye contact.
A model trained solely on single-person data can activate each face correctly, yet every identity remains blank whenever it is not speaking, an effect that looks highly unnatural in conversational scenes.

\subsection*{B.4. Future Work}
At present, AnyTalker supports only rudimentary camera motions driven by textual prompts.
Drawing inspiration from recent controllable video-generation approaches~\cite{ye2025unic, lu2024coarse, ye2025stylemaster}, we can incorporate additional conditional signals, such as camera trajectory~\cite {bai2025recammaster}.
By grafting lightweight, training-efficient modules~\cite{hu2022lora,kong2025taming,lu2025adversarial,ren2024hyper} into recent camera-trajectory-control techniques~\cite{zhang2025zero,bahmani2025ac3d}, we expect to enrich the visual storytelling of the generated videos, automatically framing and tracking the active speaker without manual intervention.

